% ============ 第1章 §1.8 Calderón-Zygmund分解与Hilbert变换的弱(1,1)性(原书 pp.54-70) ============
\section{\texorpdfstring{Calder\'on--Zygmund 分解与 Hilbert 变换的弱 $(1,1)$ 性}{Calderon--Zygmund 分解与 Hilbert 变换的弱 (1,1) 性}}\label{sec:1.8}
\index{Calderón--Zygmund 分解}
\index{弱 (1,1) 型!Hilbert 变换的}
\index{Kolmogorov 不等式}
\index{插值定理!Marcinkiewicz 的}


当 $1 < p < \infty$ 时,我们已经看到 Hilbert 变换是 $L^{p}$ 有界线性算子.那么 $p = 1$ 的情形如何呢:$L^{1}$ 有界性不成立,但称为弱 $(1,1)$ 性的结果由 A. P. Calderón 与 A. Zygmund 给出.其关键就是下面的 Calderón--Zygmund 分解.

\begin{lemma}[Calderón--Zygmund 分解]\label{lem:1.8.1}
设 $0 \leq f(x) \in L^{1}(\mathbb{R})$.则对任意 $\lambda > 0$,存在 $\mathbb{R}$ 的如下分解:

(i) $\mathbb{R} = F \cup \Omega$,$F \cap \Omega = \emptyset$;

(ii) $f(x) \leq \lambda$(a.e. $x \in F$);

(iii) $\Omega = \bigcup_{j} Q_{j}$,$Q_{j} \in \mathbb{D}$,$Q_{j} \cap Q_{k} = \emptyset$($j \neq k$),且
\begin{align*}
\text{(a)} & \lambda < \dfrac{1}{|Q_{j}|}\displaystyle\int_{Q_{j}} f(x)\dd x \leq 2\lambda, \\[4pt]
\text{(b)} & |\Omega| = \displaystyle\sum_{j}|Q_{j}| \leq \dfrac{1}{\lambda}\displaystyle\int_{\mathbb{R}} f(x)\dd x.
\end{align*}
\end{lemma}

\begin{proof}
设 $0 \leq f(x) \in L^{1}(\mathbb{R})$.固定一个满足
\begin{equation*}
\int_{\mathbb{R}} f(x)\dd x \leq 2^{k_{0}}\lambda
\end{equation*}
的整数 $k_{0}$.则对 $Q \in \mathbb{D}_{k_{0}}$,
\begin{equation}\tag{1.8.1}
\frac{1}{|Q|}\int_{Q} f(x)\dd x \leq \lambda.
\label{eq:1.8.1}\end{equation}
$\mathbb{D}_{k_{0}-1}$ 由 $\mathbb{D}_{k_{0}}$ 的二进区间二等分而得到;将 $\mathbb{D}_{k_{0}-1}$ 的二进区间 $Q$ 中满足
\begin{equation}\tag{1.8.2}
\frac{1}{|Q|}\int_{Q} f(x)\dd x > \lambda
\label{eq:1.8.2}\end{equation}
的区间称为属于第 1 群.再将不属于第 1 群的 $\mathbb{D}_{k_{0}-1}$ 的二进区间二等分所得 $\mathbb{D}_{k_{0}-2}$ 的二进区间中,满足 \eqref{eq:1.8.2} 的区间称为属于第 1 群.以下归纳地重复这一操作.如此所得属于第 1 群的二进区间的全体,因二进区间的全体是可数集,故当然可数,记为 $\{Q_{j}\}_{j}$.由作法,$Q_{j} \cap Q_{k} = \emptyset$($j \neq k$),且包含 $Q_{j}$ 的最小二进区间 $Q$ 不属于第 1 群,故满足 \eqref{eq:1.8.1}.于是
\begin{equation*}
\lambda < \frac{1}{|Q_{j}|}\int_{Q_{j}} f(x)\dd x \leq \frac{|Q|}{|Q_{j}|}\cdot\frac{1}{|Q|}\int_{Q} f(x)\dd x \leq 2\lambda.
\end{equation*}
又因 $Q_{j}$ 满足 \eqref{eq:1.8.2},
\begin{equation*}
\sum_{j}|Q_{j}| \leq \sum_{j}\frac{1}{\lambda}\int_{Q_{j}} f(x)\dd x \leq \frac{1}{\lambda}\int_{-\infty}^{\infty} f(x)\dd x.
\end{equation*}
故令 $\Omega = \bigcup_{j} Q_{j}$,$F = \Omega^{c}$,则 (i),(iii) 成立.

关于 (ii).设 $x \in F = \Omega^{c}$,则对每个 $k \leq k_{0}$,存在二进区间 $R_{k} \in \mathbb{D}_{k}$ 使 $x \in R_{k}$.于是对几乎所有的 $x \in F$,
\begin{equation*}
f(x) = \lim_{k \to -\infty}\frac{1}{|R_{k}|}\int_{R_{k}} f(y)\dd y \leq \lambda.
\end{equation*}
\end{proof}

利用这个分解,可以证明 Hilbert 变换的弱 $(1,1)$ 性.

\begin{theorem}\label{thm:1.8.2}
对 $f \in L^{1}(\mathbb{R}) \cap L^{2}(\mathbb{R})$,存在常数 $C > 0$ 使得
\begin{equation*}
\left|\{x \in \mathbb{R};\ |Hf(x)| > \lambda\}\right| \leq \frac{C}{\lambda}\| f \|_{L^{1}} \quad (\lambda > 0).
\end{equation*}
\end{theorem}

\begin{proof}
设 $f \in L^{1}(\mathbb{R}) \cap L^{2}(\mathbb{R})$,$\lambda > 0$.对 $|f(x)|$ 应用 Calderón--Zygmund 分解的引理~\ref{lem:1.8.1},定义
\begin{align*}
g(x) &= \begin{cases} f(x) & (x \in F) \\ \dfrac{1}{|Q_{j}|}\displaystyle\int_{Q_{j}} f(x)\dd x & (x \in Q_{j},\ j = 1, 2, \dots) \end{cases} \\[10pt]
f_{j}(x) &= \left(f(x) - \dfrac{1}{|Q_{j}|}\displaystyle\int_{Q_{j}} f(y)\dd y\right)\chi_{Q_{j}} \\[8pt]
b(x) &= \displaystyle\sum_{j} f_{j}(x).
\end{align*}
则 $f(x) = g(x) + b(x)$,且
\begin{equation*}
\begin{aligned}
|g(x)| &\leq 2\lambda \quad (\text{a.e. } x \in \mathbb{R}), \\
\int_{\mathbb{R}}|g(x)|\dd x &\leq \int_{\mathbb{R}}|f(x)|\dd x.
\end{aligned}
\end{equation*}
由此 $g \in L^{1}(\mathbb{R}) \cap L^{\infty}(\mathbb{R})$,故亦有 $g \in L^{1}(\mathbb{R}) \cap L^{2}(\mathbb{R})$.另外
\begin{equation*}
\int_{Q_{j}} f_{j}(x)\dd x = 0.
\end{equation*}
其次,设 $Q_{j}$ 的中心为 $x_{j}$,
\begin{equation*}
Q_{j}^{*} = 2Q_{j} = Q(x_{j}, 2\ell(Q_{j})), \quad \Omega^{*} = \bigcup_{j} Q_{j}^{*},
\end{equation*}
则
\begin{equation}\tag{1.8.3}
|\Omega^{*}| \leq \sum_{j}|Q_{j}^{*}| = \sum_{j} 2|Q_{j}| = 2|\Omega| \leq \frac{2}{\lambda}\| f \|_{L^{1}}.
\label{eq:1.8.3}\end{equation}
现在,由 $f(x) = g(x) + b(x)$ 得 $Hf(x) = Hg(x) + Hb(x)$,若 $|Hg(x)| \leq \lambda/2$ 且 $|Hb(x)| \leq \lambda/2$,则 $|Hf(x)| \leq |Hg(x)| + |Hb(x)| \leq \lambda$,故
\begin{align*}
\{x \in \mathbb{R};\ |Hf(x)| > \lambda,\ x \in (\Omega^{*})^{c}\}&\subset \{x \in (\Omega^{*})^{c};\ |Hg(x)| > \lambda/2\} \cup \{x \in (\Omega^{*})^{c};\ |Hb(x)| > \lambda/2\}.
 \end{align*}
令
\begin{equation*}
E_{1} = \{x \in (\Omega^{*})^{c};\ |Hg(x)| > \lambda/2\}, \quad E_{2} = \{x \in (\Omega^{*})^{c};\ |Hb(x)| > \lambda/2\}.
\end{equation*}
首先,利用 $H$ 的 $L^{2}$ 有界性(定理~\ref{thm:1.3.5}),
\begin{align*}
|E_{1}| &\leq \dfrac{1}{(\lambda/2)^{2}}\displaystyle\int_{E_{1}}|Hg(x)|^{2}\dd x \leq \dfrac{4}{\lambda^{2}}\displaystyle\int_{\mathbb{R}}|Hg(x)|^{2}\dd x \\
\quad &\leq \dfrac{4}{\lambda^{2}}\displaystyle\int_{\mathbb{R}}|g(x)|^{2}\dd x = \dfrac{4}{\lambda^{2}}\left(\displaystyle\int_{\Omega}|g(x)|^{2}\dd x + \displaystyle\int_{\Omega^{c}}|g(x)|^{2}\dd x\right) \\
\quad &= \dfrac{4}{\lambda^{2}}\left(\displaystyle\sum_{j}|Q_{j}||f_{Q_{j}}|^{2} + \lambda\displaystyle\int_{\Omega^{c}}|g(x)|\dd x\right) \quad (|g(x)| \leq \lambda,\ \text{a.e. } x \in \Omega^{c}) \\
\quad &\leq \dfrac{4}{\lambda^{2}}\left(\displaystyle\sum_{j}|Q_{j}|(2\lambda)^{2} + \lambda\displaystyle\int_{\mathbb{R}}|f(x)|\dd x\right) \\
&\quad \quad \quad \quad \quad \quad \quad \quad \quad \quad \quad \quad (|f_{Q_{j}}| \leq 2\lambda,\ g(x) = f(x)\ (x \in \Omega^{c})) \\
&= 16\displaystyle\sum_{j}|Q_{j}| + \dfrac{4}{\lambda}\| f \|_{L^{1}} \leq \dfrac{20}{\lambda}\| f \|_{L^{1}}.
\end{align*}
其次,若 $x \in (Q_{j}^{*})^{c}$,则
\begin{align*}
Hf_{j}(x) & = \displaystyle\lim_{\varepsilon \to 0}\frac{1}{\pi}\int_{|y - x| > \varepsilon}\frac{1}{x - y}f_{j}(y)\dd y \\
& = \dfrac{1}{\pi}\displaystyle\int_{Q_{j}}\frac{1}{x - y}f_{j}(y)\dd y - \dfrac{1}{\pi}\int_{Q_{j}}\frac{1}{x - x_{j}}f_{j}(y)\dd y \\
& \quad \left(\because \displaystyle\int_{Q_{j}} f_{j}(y)\dd y = 0\right) \\
& = \dfrac{1}{\pi}\displaystyle\int_{Q_{j}}\left(\frac{1}{x - y} - \frac{1}{x - x_{j}}\right)f_{j}(y)\dd y
\end{align*}
又若 $y \in Q_{j}$,$x \in (Q_{j}^{*})^{c}$,则 $|x - x_{j}| \geq \ell(Q_{j}) > 2|x_{j} - y|$,故 $|x - y| \geq |x - x_{j}| - |x_{j} - y| \geq |x - x_{j}|/2$,$|x - x_{j}|^{2} + \ell(Q_{j})^{2} \leq 2|x - x_{j}|^{2}$,从而
\begin{equation*}
\left|\frac{1}{x - y} - \frac{1}{x - x_{j}}\right| = \frac{|y - x_{j}|}{|x - y||x - x_{j}|} \leq \frac{2\ell(Q_{j})}{|x - x_{j}|^{2} + \ell(Q_{j})^{2}}.
\end{equation*}
于是若 $x \in (Q_{j}^{*})^{c}$,则
\begin{align*}
|Hb(x)| &\leq \displaystyle\sum_{j}|Hf_{j}(x)| \leq \sum_{j}\frac{1}{\pi}\int_{Q_{j}}\frac{2\ell(Q_{j})}{|x - x_{j}|^{2} + \ell(Q_{j})^{2}}|f_{j}(y)|\dd y \\
&\leq \dfrac{8\lambda}{\pi}\displaystyle\sum_{j}\frac{\ell(Q_{j})^{2}}{|x - x_{j}|^{2} + \ell(Q_{j})^{2}}.
\end{align*}
由此
\begin{align*}
|E_{2}| &\leq \dfrac{1}{\lambda/2}\displaystyle\int_{E_{2}}|Hb(x)|\dd x \leq \dfrac{2}{\lambda}\int_{(\Omega^{*})^{c}}|Hb(x)|\dd x \\
\quad &\leq \dfrac{16}{\pi}\displaystyle\int_{(\Omega^{*})^{c}}\sum_{j}\frac{\ell(Q_{j})^{2}}{|x - x_{j}|^{2} + \ell(Q_{j})^{2}}\dd x \\
\quad &\leq \dfrac{16}{\pi}\displaystyle\sum_{j}\ell(Q_{j})\int_{-\infty}^{\infty}\frac{\ell(Q_{j})}{|x - x_{j}|^{2} + \ell(Q_{j})^{2}}\dd x \\
\quad &= \dfrac{16}{\pi}\displaystyle\sum_{j}|Q_{j}|\int_{-\infty}^{\infty}\frac{1}{t^{2} + 1}\dd t \quad (t = (x - x_{j})/\ell(Q_{j})) \\
\quad &\leq \dfrac{16}{\lambda}\| f \|_{L^{1}}.
\end{align*}
因此
\begin{equation}\tag{1.8.4}
\left|\left\{x \in \mathbb{R};\ |Hf(x)| > \lambda,\ x \in (\Omega^{*})^{c}\right\}\right| \leq |E_{1}| + |E_{2}| \leq \frac{20 + 16}{\lambda}\| f \|_{L^{1}}.
\label{eq:1.8.4}\end{equation}
由 \eqref{eq:1.8.3},\eqref{eq:1.8.4},
\begin{align*}
|\{x \in \mathbb{R};\ |Hf(x)| > \lambda\}| & \leq |\Omega^{*}| + |\{x \in \mathbb{R};\ |Hf(x)| > \lambda,\ x \in (\Omega^{*})^{c}\}| \\
& \leq \dfrac{2 + 20 + 16}{\lambda}\| f \|_{L^{1}} = \dfrac{38}{\lambda}\| f \|_{L^{1}}.
\end{align*}
\end{proof}

\begin{note}\label{note:1.8.3}
把定理~\ref{thm:1.8.2} 证明开头所述 $f(x) = g(x) + b(x)$ 的分解,也常称为 Calderón--Zygmund 分解.另外,定理~\ref{thm:1.8.2} 可唯一地扩张到 $L^{1}(\mathbb{R})$ 上,使它对一切 $f \in L^{1}(\mathbb{R})$ 成立.这一点由下面的引理保证.
\end{note}

\begin{lemma}\label{lem:1.8.4}
\begin{sloppypar}
设 $(X, \mu)$ 为完备测度空间,$E \subset L^{1}(\mu)$ 且 $E$ 在 $L^{1}(\mu)$ 中稠密.若 $T$ 是从 $E$ 到 $\{\mu\text{-可测函数的全体}\}$ 的线性映射,满足
\end{sloppypar}
\begin{equation*}
\mu(\{x \in X;\ |Tf(x)| > \lambda\}) \leq \frac{C}{\lambda}\| f \|_{L^{1}(\mu)} \quad (\lambda > 0,\ f \in E),
\end{equation*}
则 $T$ 可唯一扩张为满足弱 $(1,1)$ 性的线性映射到 $L^{1}(\mu)$ 上.扩张后的映射仍记为 $T$,则

(i)
\begin{equation*}
\mu(\{x \in X;\ |Tf(x)| > \lambda\}) \leq \frac{C}{\lambda}\| f \|_{L^{1}(\mu)} \quad (\lambda > 0,\ f \in L^{1}(\mu)).
\end{equation*}

(ii) 若 $\lim_{j \to \infty}\| f_{j} - f \|_{L^{1}(\mu)} = 0$,则对任意 $\varepsilon > 0$ 有 $\lim_{j \to \infty}\mu(\{x \in X;\ |Tf_{j}(x) - Tf(x)| > \varepsilon\}) = 0$,即当 $j \to \infty$ 时 $Tf_{j}$ 依测度收敛于 $Tf$.
\end{lemma}

\begin{proof}
设 $f \in L^{1}(\mu) \setminus E$.因 $E$ 在 $L^{1}(\mu)$ 中稠密,存在 $f_{j} \in E$($j = 1, 2, \ldots$)满足
\begin{equation*}
\| f_{j} - f_{j+1} \|_{L^{1}(\mu)} \leq 2^{-2j}, \quad \lim_{j \to \infty}\| f_{j} - f \|_{L^{1}(\mu)} = 0.
\end{equation*}
对这个 $\{f_{j}\}$ 令
\begin{equation*}
F_{j} := \{x \in X;\ |Tf_{j}(x) - Tf_{j+1}(x)| > 2^{-j}\}.
\end{equation*}
则
\begin{equation*}
\mu\left(\bigcup_{j = k}^{\infty} F_{j}\right) \leq \sum_{j = k}^{\infty}\mu(F_{j}) \leq \sum_{j = k}^{\infty} C2^{j}\| f_{j} - f_{j+1} \|_{L^{1}(\mu)} \leq C\sum_{j = k}^{\infty}2^{-j} = C2^{-k+1}
\end{equation*}
故
\begin{equation*}
\mu\left(\bigcap_{k = 1}^{\infty}\bigcup_{j = k}^{\infty} F_{j}\right) = 0.
\end{equation*}
于是令
\begin{equation*}
N = \bigcap_{k = 1}^{\infty}\bigcup_{j = k}^{\infty} F_{j}.
\end{equation*}
若 $x \in N^{c}$,则存在编号 $k_{0}$ 使 $x \in \bigcap_{j = k_{0}}^{\infty}F_{j}^{c}$,故
\begin{equation*}
|Tf_{j}(x) - Tf_{j+1}(x)| \leq 2^{-j} \quad (j \geq k_{0}).
\end{equation*}
由此,由复数的完备性,数列 $\{Tf_{j}(x)\}_{j=1}^{\infty}$ 在 $j \to \infty$ 时收敛到某个数.将其值记为 $Tf(x)$.又当 $x \in N$ 时令 $Tf(x) = 0$,则 $Tf(x)$ 对所有 $x \in X$ 有定义,且 $\mu$-可测函数 $Tf_{j}$ 在几乎所有的点上收敛到 $Tf(x)$,故 $Tf$ 也是 $\mu$-可测函数.现在,若 $x \in \bigcap_{j = k}^{\infty}F_{j}^{c} \subset N^{c}$,则
\begin{align*}
|Tf_{k}(x) - Tf(x)| & = \displaystyle\lim_{j \to \infty}|Tf_{k}(x) - Tf_{j}(x)| \\
& \leq \displaystyle\lim_{j \to \infty}\sum_{\ell = k}^{j-1}|Tf_{\ell}(x) - Tf_{\ell+1}(x)| \leq \sum_{\ell = k}^{\infty}2^{-\ell} = 2^{1-k}.
\end{align*}
故
\begin{equation*}
\mu(\{x \in X;\ |Tf_{k}(x) - Tf(x)| > 2 \cdot 2^{-k}\}) \leq \mu\left(\bigcup_{j = k}^{\infty} F_{j}\right) \leq 2C2^{-k}.
\end{equation*}
其次,设 $\lambda > 0$,$g \in E$.取编号 $k$ 满足
\begin{align*}
\| f - g \|_{L^{1}(\mu)} &> \max(\lambda, 1)2^{-k+1}, \quad \lambda - 2^{-k+1} > \dfrac{\lambda}{2}, \\
\| f_{k} - f \|_{L^{1}(\mu)} &\leq \| f - g \|_{L^{1}(\mu)}.
\end{align*}
则当 $|Tf(x) - Tf_{k}(x)| \leq 2^{-k+1}$ 且 $|Tf_{k}(x) - Tg(x)| \leq \lambda - 2^{-k+1}$ 时,$|Tf(x) - Tg(x)| \leq |Tf(x) - Tf_{k}(x)| + |Tf_{k}(x) - Tg(x)| \leq \lambda$,故
\begin{align*}
& \mu(\{x \in X;\ |Tf(x) - Tg(x)| > \lambda\}) \\
& \qquad \leq \mu(\{x \in X;\ |Tf(x) - Tf_{k}(x)| > 2^{-k+1}\}) \\
& \qquad\qquad + \mu(\{x \in X;\ |Tf_{k}(x) - Tg(x)| > \lambda - 2^{-k+1}\}) \\
& \qquad \leq 2C2^{-k} + \dfrac{C}{\lambda - 2^{-k+1}}\| f_{k} - g \|_{L^{1}(\mu)} \\
& \qquad \leq \dfrac{C\| f - g \|_{L^{1}(\mu)}}{\max(\lambda, 1)} + \dfrac{2C}{\lambda}\| f_{k} - g \|_{L^{1}(\mu)} \\
& \qquad \leq \dfrac{C\| f - g \|_{L^{1}(\mu)}}{\lambda} + \dfrac{2C}{\lambda}\left(\| f - g \|_{L^{1}(\mu)} + \| f_{k} - f \|_{L^{1}(\mu)}\right)
\end{align*}
\begin{equation*}
\leq \frac{5C}{\lambda}\| f - g \|_{L^{1}(\mu)}.
\end{equation*}
也就是说,若 $g \in E$ 依 $L^{1}(\mu)$ 收敛于 $f$,则 $Tg$ 依测度收敛于 $Tf$.因此,定义在 $E$ 上的 $T$ 由稠密性可唯一扩张到 $L^{1}(\mu)$ 上.

进而,对 $0 < \varepsilon < 1$,因 $\{x \in X;\ |Tf(x)| > \lambda\} \subset \{x \in X;\ |Tf(x) - Tg(x)| > \varepsilon\lambda\} \cup \{x \in X;\ |Tg(x)| > (1 - \varepsilon)\lambda\}$,故
\begin{align*}
\mu(\{x \in X;\ |Tf(x)| > \lambda\}) &\leq \mu(\{x \in X;\ |Tf(x) - Tg(x)| > \varepsilon\lambda\})  \\
& \quad + \mu(\{x \in X;\ |Tg(x)| > (1 - \varepsilon)\lambda\}) \\
& \leq \dfrac{5C}{\varepsilon\lambda}\| f - g \|_{L^{1}(\mu)} + \dfrac{C}{(1 - \varepsilon)\lambda}\| g \|_{L^{1}(\mu)} \\
& \leq \left(\dfrac{5C}{\varepsilon\lambda} + \dfrac{C}{(1 - \varepsilon)\lambda}\right)\| f - g \|_{L^{1}(\mu)} + \dfrac{C}{(1 - \varepsilon)\lambda}\| f \|_{L^{1}(\mu)}.
\end{align*}
于是取 $g_{j} \in E$ 收敛于 $f$,令 $j \to \infty$ 后再令 $\varepsilon \to 0$,得
\begin{equation*}
\mu(\{x \in X;\ |Tf(x)| > \lambda\}) \leq \frac{C}{\lambda}\| f \|_{L^{1}(\mu)}.
\end{equation*}
$T$ 的线性性也容易确认.
\end{proof}

作为相关的话题,我们来叙述下面的 Kolmogorov 不等式.

\begin{theorem}[Kolmogorov 不等式]\label{thm:1.8.5}
设 $(X, \mu)$ 为测度空间.若存在常数 $A > 0$ 使得
\begin{equation*}
\mu(\{x \in X;\ |g(x)| > \lambda\}) \leq \frac{A}{\lambda} \quad (\lambda > 0),
\end{equation*}
则对任意 $0 < \delta < 1$ 与 $\mu(E) < \infty$ 的 $\mu$-可测集 $E$,存在常数 $C_{\delta}$ 使得
\begin{equation*}
\int_{E}|g(x)|^{\delta}\dd \mu \leq C_{\delta}A^{\delta}\mu(E)^{1-\delta}.
\end{equation*}
\end{theorem}

\begin{proof}
不妨设 $A = 1$.对 $k \in \mathbb{Z}$ 令
\begin{equation*}
E_{k} = \{x \in E;\ |g(x)| > 2^{k}\},
\end{equation*}
并令 $k_{0}$ 为使 $2^{-k_{0}} \geq \mu(E)$ 的最大整数.则 $\mu(E_{k}) \leq 2^{-k}$,且
\begin{align*}
\displaystyle\int_{E}|g(x)|^{\delta}\dd \mu &= \sum_{k = -\infty}^{\infty}\int_{E_{k} \setminus E_{k+1}}|g(x)|^{\delta}\dd \mu \\
&\leq \displaystyle\sum_{k = -\infty}^{\infty}2^{(k+1)\delta}\left(\mu(E_{k}) - \mu(E_{k+1})\right) \\
&= \displaystyle\sum_{k = -\infty}^{\infty}(2^{\delta} - 1)2^{k\delta}\mu(E_{k}) \\
&= \displaystyle\sum_{k < k_{0}}(2^{\delta} - 1)2^{k\delta}\mu(E_{k}) + \sum_{k \geq k_{0}}(2^{\delta} - 1)2^{k\delta}\mu(E_{k}) \\
&\leq \displaystyle\sum_{k < k_{0}}(2^{\delta} - 1)2^{k\delta}\mu(E) + \sum_{k \geq k_{0}}(2^{\delta} - 1)2^{k\delta}2^{-k} \\
&= (2^{\delta} - 1)\dfrac{2^{(k_{0}-1)\delta}}{1 - 2^{-\delta}}\mu(E) + (2^{\delta} - 1)\dfrac{2^{-k_{0}(1-\delta)}}{1 - 2^{\delta-1}} \\
&\leq (2^{\delta} - 1)\left(\dfrac{2^{-\delta}}{1 - 2^{-\delta}} + \dfrac{2^{1-\delta}}{1 - 2^{\delta-1}}\right)\mu(E)^{1-\delta}.
\end{align*}
\end{proof}

在本节最后注意如下.历史上,Hilbert 变换的 $L^{p}$ 有界性的证明顺序是:$L^{2}$ 有界性之后先证明弱 $(1,1)$ 性,对 $1 < p < 2$ 用 Marcinkiewicz 插值法证明 $L^{p}$ 有界性,对 $2 < p < \infty$ 则用对偶性来证明.用这个方法,Calderón 与 Zygmund 在 $n$ 维 Euclid 空间 $\mathbb{R}^{n}$ 上建立了``Calderón--Zygmund''奇异积分理论.Marcinkiewicz 插值法也是在证明 Hardy--Littlewood 极大函数的 $L^{p}$ 有界性时实质使用的方法,这里叙述其一般形式.

\begin{theorem}[Marcinkiewicz 插值定理]\label{thm:1.8.6}
设 $(X, \mu)$,$(Y, \nu)$ 为测度空间,$1 \leq p_{1} \leq q_{1} \leq \infty$,$1 \leq p_{2} \leq q_{2} \leq \infty$,且 $q_{1} \neq q_{2}$.设 $T$ 是从 $L^{p_{1}}(\mu) + L^{p_{2}}(\mu)$ 到 $\nu$-可测函数全体构成的函数空间的映射,满足

(i)
\begin{equation*}
|T(f_{1} + f_{2})(y)| \leq \kappa\left(|Tf_{1}(y)| + |Tf_{2}(y)|\right) \quad (\nu\text{-a.e. } y \in Y);
\end{equation*}

(ii)
\begin{align*}
& \nu(\{y \in Y;\ |Tf(y)| > \lambda\}) \leq \left(\dfrac{M_{1}}{\lambda}\| f \|_{L^{p_{1}}(\mu)}\right)^{q_{1}} \quad (\lambda > 0,\ f \in L^{p_{1}}(\mu)) \\[4pt]
& \left(q_{1} = \infty \text{ 时 } \| Tf \|_{L^{\infty}(\nu)} \leq M_{1}\| f \|_{L^{p_{1}}(\mu)}\right);
\end{align*}

(iii)
\begin{align*}
& \nu(\{y \in Y;\ |Tf(y)| > \lambda\}) \leq \left(\dfrac{M_{2}}{\lambda}\| f \|_{L^{p_{2}}(\mu)}\right)^{q_{2}} \quad (\lambda > 0,\ f \in L^{p_{2}}(\mu)) \\[4pt]
& \left(q_{2} = \infty \text{ 时 } \| Tf \|_{L^{\infty}(\nu)} \leq M_{2}\| f \|_{L^{p_{2}}(\mu)}\right).
\end{align*}

则对 $0 < t < 1$,对满足 $\dfrac{1}{p} = \dfrac{1 - t}{p_{1}} + \dfrac{t}{p_{2}}$,$\dfrac{1}{q} = \dfrac{1 - t}{q_{1}} + \dfrac{t}{q_{2}}$ 的 $p, q$,存在常数 $K$ 使得
\begin{equation*}
\| Tf \|_{L^{q}(\nu)} \leq KM_{1}^{1-t}M_{2}^{t}\| f \|_{L^{p}(\mu)} \quad (f \in L^{p}(\mu)).
\end{equation*}
此常数 $K$ 依赖于 $t, \kappa, p_{1}, p_{2}, q_{1}, q_{2}$;若 $0 < t_{1} < t < t_{2} < 1$,则仅依赖于 $t_{1}, t_{2}, \kappa, p_{1}, p_{2}, q_{1}, q_{2}$.
\end{theorem}

\begin{proof}
不妨设 $p_{1} \leq p_{2}$.设 $f \in L^{p}(\mu)$,且 $\dfrac{1}{p} = \dfrac{1 - t}{p_{1}} + \dfrac{t}{p_{2}}$,$\dfrac{1}{q} = \dfrac{1 - t}{q_{1}} + \dfrac{t}{q_{2}}$.

(a-1) $p_{1} < p_{2} < \infty$,$q_{1} < q_{2} < \infty$ 的情形.对 $a > 0$ 定义
\begin{equation*}
f_{1,a}(x) = \begin{cases} 0 & (|f(x)| \leq a) \\ f(x) & (|f(x)| > a), \end{cases}
\end{equation*}
\begin{equation*}
f_{2,a}(x) = \begin{cases} f(x) & (|f(x)| \leq a) \\ 0 & (|f(x)| > a). \end{cases}
\end{equation*}
因 $p_{1} < p$,$a^{p-p_{1}}|f_{1,a}(x)|^{p_{1}} \leq |f_{1,a}(x)|^{p}$,$|f_{1,a}(x)| \leq |f(x)| \in L^{p}(\mu)$,故
\begin{equation}\tag{1.8.5}
\int_{X}|f_{1,a}(x)|^{p_{1}}\dd \mu \leq a^{p_{1}-p}\int_{X}|f_{1,a}(x)|^{p}\dd \mu \leq a^{p_{1}-p}\| f \|_{L^{p}(\mu)}^{p} < \infty.
\label{eq:1.8.5}\end{equation}
又因 $p < p_{2}$,$|f_{2,a}(x)| \leq a$,$|f_{2,a}(x)| \leq |f(x)| \in L^{p}(\mu)$,故
\begin{equation}\tag{1.8.6}
\int_{X}|f_{2,a}(x)|^{p_{2}}\dd \mu \leq a^{p_{2}-p}\int_{X}|f_{2,a}(x)|^{p}\dd \mu \leq a^{p_{2}-p}\| f \|_{L^{p}(\mu)}^{p} < \infty.
\label{eq:1.8.6}\end{equation}
于是 $f = f_{1,a} + f_{2,a}$,$f_{1,a} \in L^{p_{1}}(\mu)$,$f_{2,a} \in L^{p_{2}}(\mu)$,故 $T(f) = T(f_{1,a} + f_{2,a})$ 有定义,由假设 (i),
\begin{equation}\tag{1.8.7}
|Tf(y)| \leq \kappa\left(|Tf_{1,a}(y)| + |Tf_{2,a}(y)|\right).
\label{eq:1.8.7}\end{equation}
这里令
\begin{equation*}
\begin{aligned}
m(\lambda) &= \mu(\{x \in X;\ |f(x)| > \lambda\}), \\
m_{1,a}(\lambda) &= \mu(\{x \in X;\ |f_{1,a}(x)| > \lambda\}), \\
m_{2,a}(\lambda) &= \mu(\{x \in X;\ |f_{2,a}(x)| > \lambda\}), \\
n(\lambda) &= \nu(\{y \in Y;\ |Tf(y)| > \lambda\}),
\end{aligned}
\end{equation*}
则
\begin{equation}\tag{1.8.8}
\| f \|_{L^{p}(\mu)}^{p} = \int_{X}|f(x)|^{p}\dd \mu = p\int_{0}^{\infty}\lambda^{p-1}m(\lambda)\dd \lambda,
\label{eq:1.8.8}\end{equation}
\begin{equation}\tag{1.8.9}
\begin{aligned}
\| Tf \|_{L^{q}(\nu)}^{q} &= \displaystyle\int_{Y}|Tf(y)|^{q}\dd \nu = q\int_{0}^{\infty}\lambda^{q-1}n(\lambda)\dd \lambda \\
&= (2\kappa)^{q}q\displaystyle\int_{0}^{\infty}\lambda^{q-1}n(2\kappa\lambda)\dd \lambda,
\end{aligned}
\label{eq:1.8.9}\end{equation}
且
\begin{equation}\tag{1.8.10}
m_{1,a}(s) = \begin{cases} m(a) & (0 < s < a), \\ m(s) & (s \geq a), \end{cases} \quad
m_{2,a}(s)\begin{cases} \leq m(s) & (0 < s < a), \\ = 0 & (s \geq a). \end{cases}
\label{eq:1.8.10}\end{equation}
其次,设 $\xi, A$ 为后面确定的常数,对 $\lambda > 0$ 令 $a = (\lambda/A)^{\xi}$.则由式~\eqref{eq:1.8.7},若 $|Tf_{1,a}(y)| \leq \lambda$ 且 $|Tf_{2,a}(y)| \leq \lambda$,则 $|Tf(y)| \leq 2\kappa\lambda$,故
\begin{equation*}
n(2\kappa\lambda) \leq \nu(\{y \in Y;\ |Tf_{1,a}(y)| > \lambda\}) + \nu(\{y \in Y;\ |Tf_{2,a}(y)| > \lambda\}).
\end{equation*}
因此,由假设 (ii),(iii) 与 \eqref{eq:1.8.5},\eqref{eq:1.8.6},
\begin{align*}
n(2\kappa\lambda) & \leq M_{1}^{q_{1}}\lambda^{-q_{1}}\| f_{1,a} \|_{L^{p_{1}}(\mu)}^{q_{1}} + M_{2}^{q_{2}}\lambda^{-q_{2}}\| f_{2,a} \|_{L^{p_{2}}(\mu)}^{q_{2}} \\
& \leq M_{1}^{q_{1}}\lambda^{-q_{1}}\left(a^{p_{1}-p}\| f \|_{L^{p}(\mu)}^{p}\right)^{q_{1}/p_{1}-1}\| f_{1,a} \|_{L^{p_{1}}(\mu)}^{p_{1}} \\
& \quad + M_{2}^{q_{2}}\lambda^{-q_{2}}\left(a^{p_{2}-p}\| f \|_{L^{p}(\mu)}^{p}\right)^{q_{2}/p_{2}-1}\| f_{2,a} \|_{L^{p_{2}}(\mu)}^{p_{2}}.
\end{align*}
由此与 \eqref{eq:1.8.9},
\begin{align*}
\| Tf \|_{L^{q}(\nu)}^{q} & = (2\kappa)^{q}qM_{1}^{q_{1}}\| f \|_{L^{p}(\mu)}^{p(q_{1}/p_{1}-1)} \\
& \quad \times \displaystyle\int_{0}^{\infty}\lambda^{q-q_{1}-1}a^{(p_{1}-p)(q_{1}/p_{1}-1)}\| f_{1,a} \|_{L^{p_{1}}(\mu)}^{p_{1}}\dd \lambda \\
& \quad + (2\kappa)^{q}qM_{2}^{q_{2}}\| f \|_{L^{p}(\mu)}^{p(q_{2}/p_{2}-1)} \\
& \quad \times \displaystyle\int_{0}^{\infty}\lambda^{q-q_{2}-1}a^{(p_{2}-p)(q_{2}/p_{2}-1)}\| f_{2,a} \|_{L^{p_{2}}(\mu)}^{p_{2}}\dd \lambda \\
& =: (2\kappa)^{q}qM_{1}^{q_{1}}\| f \|_{L^{p}(\mu)}^{p(q_{1}/p_{1}-1)}I_{1} + (2\kappa)^{q}qM_{2}^{q_{2}}\| f \|_{L^{p}(\mu)}^{p(q_{2}/p_{2}-1)}I_{2}.
\end{align*}
(最后一行表示把倒数第二行中的积分部分分别记为 $I_{1}, I_{2}$.)

这里令 $\xi = \dfrac{1/q_{1} - 1/q_{2}}{1/p_{1} - 1/p_{2}}\cdot\dfrac{q}{p}$.此时,
\begin{equation}\tag{1.8.11}
\xi = \frac{q - q_{1}}{p - p_{1}}\cdot\frac{p_{1}}{q_{1}} = \frac{q - q_{2}}{p - p_{2}}\cdot\frac{p_{2}}{q_{2}},
\label{eq:1.8.11}\end{equation}
\begin{equation}\tag{1.8.12}
\frac{q - q_{1}}{\xi} + (p_{1} - p)\left(\frac{q_{1}}{p_{1}} - 1\right) + p_{1} = p,
\label{eq:1.8.12}\end{equation}
\begin{equation}\tag{1.8.13}
\frac{q - q_{2}}{\xi} + (p_{2} - p)\left(\frac{q_{2}}{p_{2}} - 1\right) + p_{2} = p
\label{eq:1.8.13}\end{equation}
成立(实际上,正是为了满足上式才这样确定 $\xi$).因令 $a = (\lambda/A)^{\xi}$,故利用 \eqref{eq:1.8.10} 得
\begin{align*}
I_{1}&= A^{-\xi(p_{1}-p)(q_{1}/p_{1}-1)}\displaystyle\int_{0}^{\infty}\lambda^{q-q_{1}-1+\xi(p_{1}-p)(q_{1}/p_{1}-1)}\left(\int_{0}^{\infty}p_{1}s^{p_{1}-1}m_{1,a}(s)\dd s\right)\dd \lambda  \\
&= A^{-\xi(p_{1}-p)(q_{1}/p_{1}-1)}\displaystyle\int_{0}^{\infty}\lambda^{q-q_{1}-1+\xi(p_{1}-p)(q_{1}/p_{1}-1)}\left(\int_{(\lambda/A)^{\xi}}^{\infty}p_{1}s^{p_{1}-1}m(s)\dd s\right)\dd \lambda \\
&\quad + A^{-\xi(p_{1}-p)(q_{1}/p_{1}-1)}\displaystyle\int_{0}^{\infty}\lambda^{q-q_{1}-1+\xi(p_{1}-p)(q_{1}/p_{1}-1)}\left(\frac{\lambda}{A}\right)^{\xi p_{1}}m\left(\left(\frac{\lambda}{A}\right)^{\xi}\right)\dd \lambda \\
&= A^{-\xi(p_{1}-p)(q_{1}/p_{1}-1)}\displaystyle\int_{0}^{\infty}p_{1}s^{p_{1}-1}m(s)\left(\int_{0}^{As^{1/\xi}}\lambda^{q-q_{1}-1+\xi(p_{1}-p)(q_{1}/p_{1}-1)}\dd \lambda\right)\dd s \\
&\quad + A^{-\xi(p_{1}-p)(q_{1}/p_{1}-1)}\displaystyle\int_{0}^{\infty}\lambda^{q-q_{1}-1+\xi(p_{1}-p)(q_{1}/p_{1}-1)}\left(\frac{\lambda}{A}\right)^{\xi p_{1}}m\left(\left(\frac{\lambda}{A}\right)^{\xi}\right)\dd \lambda \\
&= \frac{p_{1}A^{q-q_{1}}}{q - q_{1} + \xi(p_{1}-p)(q_{1}/p_{1}-1)}\displaystyle\int_{0}^{\infty}s^{p_{1}-1+(q-q_{1})/\xi+(p_{1}-p)(q_{1}/p_{1}-1)}m(s)\dd s
\end{align*}
\begin{equation*}
+ \frac{A^{q-q_{1}}}{\xi}\int_{0}^{\infty}s^{p_{1}-1+(q-q_{1})/\xi+(p_{1}-p)(q_{1}/p_{1}-1)}m(s)\dd s.
\end{equation*}
(在由倒数第二个式子的第一个积分项变为最后一个式子的过程中,除了对倒数第二,第三个式子的第一项使用 Fubini 定理外,还利用了 \eqref{eq:1.8.12},\eqref{eq:1.8.11} 推出的 $q - q_{1} - 1 + \xi(p_{1}-p)\left(\frac{q_{1}}{p_{1}}-1\right) = (q-q_{1})p_{1}/q_{1} - 1 > -1$.最后一步中作了变量替换 $\lambda = As^{1/\xi}$.)

对 $I_{2}$ 作同样的考察,
\begin{align*}
I_{2}&= A^{-\xi(p_{2}-p)(q_{2}/p_{2}-1)}\displaystyle\int_{0}^{\infty}\lambda^{q-q_{2}-1+\xi(p_{2}-p)(q_{2}/p_{2}-1)}\left(\int_{0}^{\infty}p_{2}s^{p_{2}-1}m_{2,a}(s)\dd s\right)\dd \lambda  \\
&\leq A^{-\xi(p_{2}-p)(q_{2}/p_{2}-1)}\displaystyle\int_{0}^{\infty}\lambda^{q-q_{2}-1+\xi(p_{2}-p)(q_{2}/p_{2}-1)}\left(\int_{0}^{(\lambda/A)^{\xi}}p_{2}s^{p_{2}-1}m(s)\dd s\right)\dd \lambda \\
&= A^{-\xi(p_{2}-p)(q_{2}/p_{2}-1)}\displaystyle\int_{0}^{\infty}p_{2}s^{p_{2}-1}m(s)\left(\int_{As^{1/\xi}}^{\infty}\lambda^{q-q_{2}-1+\xi(p_{2}-p)(q_{2}/p_{2}-1)}\dd \lambda\right)\dd s \\
&= \frac{-p_{2}A^{q-q_{2}}}{q - q_{2} + \xi(p_{2}-p)(q_{2}/p_{2}-1)}\displaystyle\int_{0}^{\infty}s^{p_{2}-1+(q-q_{2})/\xi+(p_{2}-p)(q_{2}/p_{2}-1)}m(s)\dd s
\end{align*}
因此,利用 $\xi$ 的定义,\eqref{eq:1.8.11},\eqref{eq:1.8.12},\eqref{eq:1.8.13} 整理得
\begin{align*}
I_{1} &= \dfrac{1/p_{1} - 1/p_{2}}{1/q_{1} - 1/q_{2}}\dfrac{pA^{q-q_{1}}}{q(p - p_{1})}\| f \|_{L^{p}(\mu)}^{p}, \\
I_{2} &\leq \dfrac{1/p_{1} - 1/p_{2}}{1/q_{1} - 1/q_{2}}\dfrac{p_{2}A^{q-q_{2}}}{q(p_{2} - p)}\| f \|_{L^{p}(\mu)}^{p}.
\end{align*}
由 $I_{1}, I_{2}$ 的估计式,
\begin{align*}
\| Tf \|_{L^{q}(\nu)}^{q} & \leq (2\kappa)^{q}M_{1}^{q_{1}}A^{q-q_{1}}\dfrac{1/p_{1} - 1/p_{2}}{1/q_{1} - 1/q_{2}}\dfrac{p}{p - p_{1}}\| f \|_{L^{p}(\mu)}^{pq_{1}/p_{1}} \\
& \quad + (2\kappa)^{q}M_{2}^{q_{2}}A^{q-q_{2}}\dfrac{1/p_{1} - 1/p_{2}}{1/q_{1} - 1/q_{2}}\dfrac{p_{2}}{p_{2} - p}\| f \|_{L^{p}(\mu)}^{pq_{2}/p_{2}}.
\end{align*}
于是,选取 $A$ 使
\begin{equation*}
M_{1}^{q_{1}}A^{q-q_{1}}\| f \|_{L^{p}(\mu)}^{pq_{1}/p_{1}} = M_{2}^{q_{2}}A^{q-q_{2}}\| f \|_{L^{p}(\mu)}^{pq_{2}/p_{2}},
\end{equation*}
即令
\begin{equation*}
A = M_{1}^{q_{1}/(q_{1}-q_{2})}M_{2}^{q_{2}/(q_{2}-q_{1})}\| f \|_{L^{p}(\mu)}^{p(q_{2}/p_{2}-q_{1}/p_{1})/(q_{2}-q_{1})},
\end{equation*}
并注意到
\begin{equation*}
\frac{q - q_{1}}{q_{2} - q_{1}}q_{2} = qt, \qquad
\frac{pq_{1}}{p_{1}} + \frac{q - q_{1}}{q_{2} - q_{1}}\left(\frac{q_{2}}{p_{2}} - \frac{q_{1}}{p_{1}}\right)p = q,
\end{equation*}
则
\begin{equation*}
\| Tf \|_{L^{q}(\nu)}^{q} \leq (2\kappa)^{q}\frac{1/p_{1} - 1/p_{2}}{1/q_{1} - 1/q_{2}}\left(\frac{p}{p - p_{1}} + \frac{p_{2}}{p_{2} - p}\right)M_{1}^{q(1-t)}M_{2}^{qt}\| f \|_{L^{p}(\mu)}^{q}
\end{equation*}
因此
\begin{equation*}
\| Tf \|_{L^{q}(\nu)} \leq (2\kappa)\left(\frac{1/p_{1} - 1/p_{2}}{1/q_{1} - 1/q_{2}}\left(\frac{p}{p - p_{1}} + \frac{p_{2}}{p_{2} - p}\right)\right)^{1/q}M_{1}^{1-t}M_{2}^{t}\| f \|_{L^{p}(\mu)}
\end{equation*}

(a-2) $p_{1} < p_{2} < \infty$,$q_{2} < q_{1} < \infty$ 的情形.与 (a-1) 的情形几乎相同,但令 $\xi = \dfrac{1/q_{1} - 1/q_{2}}{1/p_{1} - 1/p_{2}}\cdot\dfrac{q}{p}$ 时,$\xi$ 为负这一差别会出现.从那里开始,
\begin{align*}
I_{1}&= A^{-\xi(p_{1}-p)(q_{1}/p_{1}-1)}\displaystyle\int_{0}^{\infty}p_{1}s^{p_{1}-1}m(s)\left(\int_{As^{1/\xi}}^{\infty}\lambda^{q-q_{1}-1+\xi(p_{1}-p)(q_{1}/p_{1}-1)}\dd \lambda\right)\dd s  \\
&\quad + A^{-\xi(p_{1}-p)(q_{1}/p_{1}-1)}\displaystyle\int_{0}^{\infty}\lambda^{q-q_{1}-1+\xi(p_{1}-p)(q_{1}/p_{1}-1)}\left(\frac{\lambda}{A}\right)^{\xi p_{1}}m\left(\left(\frac{\lambda}{A}\right)^{\xi}\right)\dd \lambda \\
&= \frac{-p_{1}A^{q-q_{1}}}{q - q_{1} + \xi(p_{1}-p)(q_{1}/p_{1}-1)}\displaystyle\int_{0}^{\infty}s^{p_{1}-1+(q-q_{1})/\xi+(p_{1}-p)(q_{1}/p_{1}-1)}m(s)\dd s \\
&\quad - \frac{A^{q-q_{1}}}{\xi}\displaystyle\int_{0}^{\infty}s^{p_{1}-1+(q-q_{1})/\xi+(p_{1}-p)(q_{1}/p_{1}-1)}m(s)\dd s \\
&= \dfrac{1/p_{1} - 1/p_{2}}{1/q_{2} - 1/q_{1}}\dfrac{pA^{q-q_{1}}}{q(p - p_{1})}\| f \|_{L^{p}(\mu)}^{p}. \\
&I_{2} \\
&\leq A^{-\xi(p_{2}-p)(q_{2}/p_{2}-1)}\displaystyle\int_{0}^{\infty}p_{2}s^{p_{2}-1}m(s)\left(\int_{0}^{As^{1/\xi}}\lambda^{q-q_{2}-1+\xi(p_{2}-p)(q_{2}/p_{2}-1)}\dd \lambda\right)\dd s \\
&= \dfrac{p_{2}A^{q-q_{2}}}{q - q_{2} + \xi(p_{2}-p)(q_{2}/p_{2}-1)}\displaystyle\int_{0}^{\infty}s^{p_{2}-1+(q-q_{2})/\xi+(p_{2}-p)(q_{2}/p_{2}-1)}m(s)\dd s
\end{align*}
\begin{equation*}
= \frac{1/p_{1} - 1/p_{2}}{1/q_{2} - 1/q_{1}}\frac{p_{2}A^{q-q_{2}}}{q(p_{2} - p)}\| f \|_{L^{p}(\mu)}^{p}
\end{equation*}
以下与 (a-1) 同样选取 $A$,得
\begin{equation*}
\| Tf \|_{L^{q}(\nu)} \leq (2\kappa)\left(\frac{1/p_{1} - 1/p_{2}}{1/q_{2} - 1/q_{1}}\left(\frac{p}{p - p_{1}} + \frac{p_{2}}{p_{2} - p}\right)\right)^{1/q}M_{1}^{1-t}M_{2}^{t}\| f \|_{L^{p}(\mu)}.
\end{equation*}

(a-3) $p_{1} < p_{2} < \infty$,$q_{1} < q_{2} = \infty$ 的情形.在 (a-1) 中令 $a = (\lambda/A)^{\xi}$,$\xi = p_{2}/(p_{2} - p)$,$A = M_{2}\| f \|_{L^{p}(\mu)}^{p/p_{2}}$,则由假设 (iii),
\begin{equation*}
\| Tf_{2,a}(y) \|_{L^{\infty}(\nu)} \leq M_{2}\| f_{2,a} \|_{L^{p_{2}}(\mu)} \leq M_{2}\left(a^{p_{2}-p}\| f \|_{L^{p}(\mu)}^{p}\right)^{1/p_{2}} \leq \lambda,
\end{equation*}
故 $\nu(\{y \in Y;\ |Tf_{2,a}(y)| > \lambda\}) = 0$.因此 $n(2\kappa\lambda) \leq \nu(\{y \in Y;\ |Tf_{1,a}(y)| > \lambda\})$,从而
\begin{align*}
\| Tf \|_{L^{q}(\nu)}^{q} & \leq (2\kappa)^{q}M_{1}^{q_{1}}A^{q-q_{1}}\dfrac{1/p_{1} - 1/p_{2}}{1/q_{1}}\dfrac{p}{p - p_{1}}\| f \|_{L^{p}(\mu)}^{pq_{1}/p_{1}} \\
& = (2\kappa)^{q}M_{1}^{q_{1}}M_{2}^{q-q_{1}}\dfrac{1/p_{1} - 1/p_{2}}{1/q_{1}}\dfrac{p}{p - p_{1}}\| f \|_{L^{p}(\mu)}^{p(q_{1}/p_{1} + (q-q_{1})/p_{2})}.
\end{align*}
注意 $1/q = (1-t)/q_{1}$,$q_{1}/(qp_{1}) + (1 - q_{1}/q)/p_{2} = (1-t)/p_{1} + t/p_{2} = 1/p$,计算得
\begin{equation*}
\| Tf \|_{L^{q}(\nu)} \leq 2\kappa M_{1}^{1-t}M_{2}^{t}\left(\frac{1/p_{1} - 1/p_{2}}{1/q_{1}}\frac{p}{p - p_{1}}\right)^{1/q}\| f \|_{L^{p}(\mu)}.
\end{equation*}

(a-4) $p_{1} < p_{2} < \infty$,$q_{2} < q_{1} = \infty$ 的情形.在 (a-1) 中令 $a = (\lambda/A)^{\xi}$,$\xi = p_{1}/(p_{1} - p)$,$A = M_{1}\| f \|_{L^{p}(\mu)}^{p/p_{1}}$,则由假设 (ii),
\begin{equation*}
\| Tf_{1,a}(y) \|_{L^{\infty}(\nu)} \leq M_{1}\| f_{1,a} \|_{L^{p_{1}}(\mu)} \leq M_{1}\left(a^{p_{1}-p}\| f \|_{L^{p}(\mu)}^{p}\right)^{1/p_{1}} = \lambda,
\end{equation*}
故 $\nu(\{y \in Y;\ |Tf_{1,a}(y)| > \lambda\}) = 0$.因此 $n(2\kappa\lambda) \leq \nu(\{y \in Y;\ |Tf_{2,a}(y)| > \lambda\})$,从而
\begin{align*}
\| Tf \|_{L^{q}(\nu)}^{q} & \leq (2\kappa)^{q}M_{2}^{q_{2}}A^{q-q_{2}}\dfrac{1/p_{1} - 1/p_{2}}{1/q_{2}}\dfrac{p_{2}}{p_{2} - p}\| f \|_{L^{p}(\mu)}^{pq_{2}/p_{2}} \\
& = (2\kappa)^{q}M_{1}^{q-q_{2}}M_{2}^{q_{2}}\dfrac{1/p_{1} - 1/p_{2}}{1/q_{2}}\dfrac{p}{p_{2} - p}\| f \|_{L^{p}(\mu)}^{p(q-q_{2})/p_{1} + pq_{2}/p_{2}}.
\end{align*}
注意 $1/q = t/q_{2}$,$(1 - q_{2}/q)/p_{1} + q_{2}/(qp_{2}) = (1-t)/p_{1} + t/p_{2} = 1/p$,计算得
\begin{equation*}
\| Tf \|_{L^{q}(\nu)} \leq 2\kappa M_{1}^{1-t}M_{2}^{t}\left(\frac{1/p_{1} - 1/p_{2}}{1/q_{2}}\frac{p}{p_{2} - p}\right)^{1/q}\| f \|_{L^{p}(\mu)}.
\end{equation*}

(b) $p_{1} < p_{2} = \infty$,$q_{1} < q_{2} = \infty$ 的情形.在 (a-1) 中令 $a = (\lambda/A)^{\xi}$,$\xi = 1$,$A = M_{2}$,则由假设 (iii),
\begin{equation*}
\| Tf_{2,a}(y) \|_{L^{\infty}(\nu)} \leq M_{2}\| f_{2,a} \|_{L^{\infty}(\mu)} \leq M_{2}a = \lambda,
\end{equation*}
以下与 (a-3) 同样进行,得
\begin{equation*}
\| Tf \|_{L^{q}(\nu)} \leq 2\kappa M_{1}^{1-t}M_{2}^{t}\left(\frac{q_{1}}{p_{1}}\frac{p}{p - p_{1}}\right)^{1/q}\| f \|_{L^{p}(\mu)}.
\end{equation*}

(c-1) $p_{1} = p_{2} < \infty$,$q_{1} \neq q_{2}$,$q_{1}, q_{2} < \infty$ 的情形.此时不妨设 $p_{1} = p_{2} = p$,$q_{1} < q_{2}$.由假设 (ii),(iii),
\begin{equation*}
n(\lambda) \leq \left(\frac{M_{1}}{\lambda}\| f \|_{L^{p}(\mu)}\right)^{q_{1}}, \quad n(\lambda) \leq \left(\frac{M_{2}}{\lambda}\| f \|_{L^{p}(\mu)}\right)^{q_{2}},
\end{equation*}
故对 $A > 0$,
\begin{align*}
\| Tf \|_{L^{q}(\nu)}^{q} & = \displaystyle\int_{0}^{A}q\lambda^{q-1}n(\lambda)\dd \lambda + \int_{A}^{\infty}q\lambda^{q-1}n(\lambda)\dd \lambda \\
& \leq \displaystyle\int_{0}^{A}q\lambda^{q-q_{1}-1}n(\lambda)\dd \lambda\left(M_{1}\| f \|_{L^{p}(\mu)}\right)^{q_{1}} \\
& \quad + \displaystyle\int_{A}^{\infty}q\lambda^{q-q_{2}-1}n(\lambda)\dd \lambda\left(M_{2}\| f \|_{L^{p}(\mu)}\right)^{q_{2}} \\
& = M_{1}^{q_{1}}\dfrac{q}{q - q_{1}}A^{q-q_{1}}\| f \|_{L^{p}(\mu)}^{q_{1}} + M_{2}^{q_{2}}\dfrac{q}{q_{2} - q}A^{q-q_{2}}\| f \|_{L^{p}(\mu)}^{q_{2}}.
\end{align*}
与 (a-1) 的情形一样,令 $A = M_{1}^{q_{1}/(q-q_{1})}M_{2}^{-q_{2}/(q-q_{1})}\| f \|_{L^{p}(\mu)}^{q_{1}-q_{2}}$,得
\begin{equation*}
\| Tf \|_{L^{q}(\nu)} \leq \left(\frac{1}{q - q_{1}} + \frac{1}{q_{2} - q}\right)^{1/q}M_{1}^{1-t}M_{2}^{t}\| f \|_{L^{p}(\mu)}.
\end{equation*}

(c-2) $p_{1} = p_{2} < \infty$,$q_{1}$ 或 $q_{2} = \infty$ 的情形.此时不妨设 $p_{1} = p_{2} = p$,$q_{1} < q_{2} = \infty$.由假设 (ii),(iii),
\begin{equation*}
n(\lambda) \leq \left(\frac{M_{1}}{\lambda}\| f \|_{L^{p}(\mu)}\right)^{q_{1}}, \quad \| Tf(y) \|_{L^{\infty}(\nu)} \leq M_{2}\| f \|_{L^{p}(\mu)},
\end{equation*}
故当 $\lambda > M_{2}\| f \|_{L^{p}(\mu)}$ 时 $\{y \in Y;\ |Tf(y)| > \lambda\} = \emptyset$.因此
\begin{align*}
\| Tf \|_{L^{q}(\nu)}^{q} &= \displaystyle\int_{0}^{\infty}q\lambda^{q-1}n(\lambda)\dd \lambda \\
&= \displaystyle\int_{0}^{M_{2}\| f \|_{L^{p}(\mu)}}q\lambda^{q-q_{1}-1}n(\lambda)\dd \lambda\left(M_{1}\| f \|_{L^{p}(\mu)}\right)^{q_{1}} \\
&= q\left(M_{1}\| f \|_{L^{p}(\mu)}\right)^{q_{1}}\dfrac{\left(M_{2}\| f \|_{L^{p}(\mu)}\right)^{q-q_{1}}}{q - q_{1}} = \dfrac{qM_{1}^{q_{1}}M_{2}^{q-q_{1}}}{q - q_{1}}\| f \|_{L^{p}(\mu)}^{q}.
\end{align*}
注意 $1/q = (1-t)/q_{1}$,计算得
\begin{equation*}
\| Tf \|_{L^{q}(\nu)} \leq \left(\frac{q}{q - q_{1}}\right)^{1/q}M_{1}^{1-t}M_{2}^{t}\| f \|_{L^{p}(\mu)}.
\end{equation*}

(c-3) $p_{1} = p_{2} = \infty$,$q_{1} < q_{2} = \infty$ 的情形.此时与 (c-2) 中 $p_{1} = p_{2} = p$,$q_{1} < q_{2} = \infty$ 情形的证明相同.
\end{proof}
